\chapter{Nested Samplers}
\label{ch:nested}

\newthought{Nested sampling} computes the Bayesian \emph{evidence} (the marginal likelihood) and
produces posterior samples as a by-product. It is the tool of choice when you need model
comparison or robust handling of multiple modes. \Jarvis\ offers two implementations,
\sampler{Dynesty} and \sampler{MultiNest}, which share the same configuration.

\section{The shared idea}
\label{sec:nested-idea}

Nested sampling\cite{Skilling:2006gxv} keeps a fixed set of \yamlkey{nlive} ``live points'' drawn
from the prior. At each step it removes the live point with the \emph{lowest} likelihood---recording
it as a posterior sample weighted by the shrinking prior volume it represented---and replaces it
with a fresh point of \emph{higher} likelihood. Repeating this marches the live set inward through
nested shells of increasing likelihood, accumulating the evidence integral along the way. Because
the live set can split across separated peaks, the method handles multi-modal problems naturally.

\begin{description}[style=nextline, leftmargin=1.2em]
  \item[More live points] $\Rightarrow$ more accurate evidence and finer posterior, but more
        evaluations. \yamlkey{nlive} is the main accuracy/cost dial.
\end{description}

\section{Shared format}
\label{sec:nested-format}

Both samplers use the same \yamlkey{Bounds} keys:

\begin{description}[style=nextline, leftmargin=1.2em]
  \item[\yamlkey{nlive}] the number of live points (accuracy vs.\ cost).
  \item[\yamlkey{rseed}] random seed for reproducibility.
  \item[\yamlkey{run\_nested}] run controls, such as \yamlkey{maxcall} (a cap on likelihood calls)
        and \yamlkey{print\_progress} (live progress output).
\end{description}

\begin{yamlwin}
Sampling:
  Method: "Dynesty"
  Variables: [ ... ]
  Bounds:
    nlive: 500
    rseed: 42
    run_nested:
      maxcall: 200000
      print_progress: true
  LogLikelihood:
    - {name: "LogL", expression: "..."}
\end{yamlwin}

\section{Dynesty}
\label{sec:s-dynesty}

\paragraph{How it differs} \sampler{Dynesty}\cite{Speagle:2019ivv} is a pure-Python dynamic
nested sampler: it can allocate live points adaptively during the run, putting more effort into the
posterior (or the evidence) depending on what you care about. It is a flexible default for
moderate dimensions.

\begin{jnote}
\sampler{Dynesty} ships in the optional \code{legacy} dependency set. If it is not already present,
install it before selecting this method (see Chapter~\ref{ch:installation}).
\end{jnote}

\section{MultiNest}
\label{sec:s-multinest}

\paragraph{How it differs} \sampler{MultiNest}\cite{Feroz:2008xx} draws new live points using
\emph{ellipsoidal decomposition}: it wraps the current live set in one or more ellipsoids and
samples within them. This is particularly effective for strongly multi-modal posteriors and
curving degeneracies. It relies on the external MultiNest library being installed.

\section{Reading the output}
\label{sec:nested-output}

Beyond the usual sample table, the key things to check are the reported evidence and its stability,
and the live-point behaviour over the run. The evidence is what enables model comparison: the
difference in log-evidence between two models is the log Bayes factor.
